\originalsection{译者术语表}\label{cat-glossary}
下表为译者附加，不属于原论文正文.
\Needspace{8\baselineskip}
\begin{longtable}{p{.30\textwidth}p{.22\textwidth}p{.40\textwidth}}\toprule
原术语&译法&说明\\\midrule
\endfirsthead
\toprule 原术语&译法&说明\\\midrule
\endhead
Compute as Teacher (CaT)&计算即教师&框架名\\
reference-free supervision&无参考监督&无外部人工参考，不表示无任何先验\\
rollout&采样轨迹&模型针对提示生成的一次完整回答或推理\\
reference estimation&参考估计&从轨迹集合构建监督估计\\
pseudo-reference&伪参考&模型生成目标，不等于外部真值\\
reward derivation&奖励推导&把伪参考转为每条轨迹奖励\\
synthesis&合成&综合轨迹片段生成新回答\\
selection&选择&从已有轨迹中选一个\\
anchor&锚点&冻结初始策略\\
self-proposed rubrics&自提出评分细则&由模型基于伪参考生成的二元标准\\
criterion&标准&细则中的单条可判定要求\\
judge&判别器&判断是否满足标准的模型\\
verifiable domain&可验证领域&原文指可程序检查答案的领域\\
non-verifiable domain&不可验证领域&原文指定性答案且缺乏程序检查器的领域\\
best-of-N / Self-BoN&best-of-N／Self-BoN&从N个样本挑最好；Self指策略自行评选\\
mutual predictability&互可预测性&在其他回答条件下评价当前回答概率\\
group-relative advantage&组相对优势&相对组均值并以组标准差归一化\\
amortize&摊入／摊销&通过训练让后续单次生成复用聚合收益\\
reward hacking&奖励投机&提高奖励却偏离任务目的\\
entropy collapse&熵坍缩&策略输出多样性降低\\
oracle baseline&理想基线&使用专家标注的对照\\
standard error&标准误&跨测试样本的均值不确定性指标\\
Welch t-test&Welch t检验&原文采用的组间显著性检验\\
surface-form bias&表面形式偏差&对措辞、长度或格式的偏好\\
\bottomrule\end{longtable}
